% !TEX root = ../tecnologia-en-marcha.tex
% !TEX encoding = UTF-8 Unicode
% Texto: borrador de Laura Segura, con acentos restaurados y citas enlazadas.

\section{Introducción}

Los sensores ambientales instalados en viviendas han sido utilizados para
reconocer actividades de la vida diaria sin recurrir a cámaras ni micrófonos
\cite{kim2022inhome,requena2024independence}. En particular, los sensores de infrarrojo pasivo, junto con sensores magnéticos de
puerta y sensores de temperatura, permiten registrar patrones de movimiento
compatibles con tareas como dormir, cocinar, comer o desplazarse dentro del
hogar. Esta clase de sistemas resulta relevante para el monitoreo no intrusivo
de personas adultas mayores, pues puede aportar información de apoyo a
cuidadores y servicios de acompañamiento \cite{park2025aging}. Sin embargo,
cuando las inferencias automatizadas se usan sobre poblaciones vulnerables, el
desempeño técnico del clasificador no es suficiente para justificar su uso.
% Las dos afirmaciones siguientes resumen la lectura que hace PIA02 (1.9 y
% 2.3.2) de esas revisiones.
Las revisiones recientes señalan, además, que la mayoría de los estudios del
campo se concentran en demostrar el desempeño técnico de los modelos y no en
evaluar su transparencia, su equidad o su gobernanza
\cite{park2025aging,habibovic2025aging}.

La Estrategia Nacional de Inteligencia Artificial de Costa Rica plantea
principios rectores para el desarrollo y uso responsable de sistemas de IA,
entre ellos transparencia y explicabilidad, equidad y no discriminación, y
responsabilidad \cite{micitt2024enia}. Estos principios son coherentes con
instrumentos internacionales como el Reglamento Europeo de Inteligencia
Artificial, que establece obligaciones de transparencia, gobernanza de datos y
conservación de registros para sistemas de alto riesgo \cite{ue2024aiact}, y con
el NIST AI Risk Management Framework, que organiza la gestión de riesgos en
funciones de mapear, medir, gobernar y gestionar \cite{nist2023airmf}. No
obstante, estos marcos suelen expresarse en un nivel normativo o metodológico
que no siempre indica cómo producir evidencia técnica verificable sobre un
sistema específico.

El problema abordado en este estudio es la brecha entre principios regulatorios
de alto nivel y mecanismos computacionales reproducibles para auditarlos en
sistemas de IA basados en sensores domiciliarios. En este dominio, la brecha es
particularmente relevante porque los datos describen rutinas domésticas
sensibles, los metadatos demográficos disponibles son limitados y los errores
pueden afectar de forma desigual a personas con patrones de movilidad o rutina
menos representados.

El objetivo de la investigación fue diseñar, implementar y verificar un marco
operativo que traduzca tres principios de la Estrategia Nacional de IA de Costa
Rica en requerimientos técnicos, pruebas computacionales y artefactos auditables
aplicables a un clasificador de actividades basado en sensores PIR. El artículo
resume la arquitectura del marco, la metodología de implementación y los
resultados obtenidos en dos aplicaciones piloto sobre datos públicos del
proyecto CASAS.
